\section{Introduction}
\label{sec:introduction}

Named Entity Recognition (NER) extracts entities such as persons, organisations and dates from raw text. It is a foundational step in most information extraction pipelines. Standard NER systems are closed-set. They are trained once on a fixed inventory of types, and can only ever predict those types. Real-world applications rarely afford such a static inventory, and legal text is an extreme case. Entity inventories there are open and shifting, with new statutes, case citations, parties, courts, contractual instruments and jurisdictions. Documents are long and formulaic, annotation requires expert effort, and extraction errors carry material consequences. Re-annotating and retraining a tagger for every new type or sub-domain is rarely practical.

\emph{Open-world} NER targets precisely this regime. It recognises and types entities whose categories were not fixed at training time. Notable systems reach it by different routes. They score spans against type names \citep{zaratiana2024gliner}, tag generatively \citep{ding2024gner}, cluster a dedicated entity encoder into dynamically named types \citep{genest2025owner}, or prompt instruction-tuned models at far higher inference compute \citep{zhou2024universalner}. All of them are trained and evaluated on general-domain text. Their legal uses then adapt them per corpus, schema and language, either by fine-tuning a tagger or by prompting a large model anew. Few legal deployments can pay that price for every schema, sub-domain and jurisdiction. We take the opposite route, with one open-world pipeline whose components are never fitted to a legal corpus. We measure it component by component, to locate the bottleneck and show where the next annotation or compute budget should go.

We do this with OPENER \rev{(Open Partitioning Embedding for Named Entity Recognition)}, a frugal open-world pipeline \rev{assembled from} ready-to-use pretrained components. A frozen open-world detector proposes spans. A light contrastively fine-tuned, Matryoshka-truncatable embedder then represents them. That embedder is fine-tuned on general-domain entity spans, and never on a legal schema. This is a design choice rather than a shortcut. An open-world space must stay unspecialised, or it collapses into the closed-set taggers we are trying to avoid. The pipeline also runs one to two orders of magnitude cheaper in latency and energy than trained-encoder and generative baselines. Legal practice is rarely GPU-rich, and needs exactly that. We evaluate its two operating points. OPENER-ZS is a label-free head, built on label-name prototypes, transductive refinement and detector fusion. OPENER-Sup is a balanced linear probe, fit on the legal target labels. Our contributions are:
\begin{itemize}
\item \textbf{A frugal pipeline for legal NER.} We assemble an open-world pipeline from pretrained components, with \emph{no} legal-domain training of any of them, and we test both of its typing heads.
\item \textbf{A legal open-world benchmark.} Four legal NER corpora in three languages, English (SEC/EDGAR filings and Indian court judgements), Portuguese and German (Brazilian and federal court decisions). Eight systems are compared end to end under one protocol, both OPENER heads, the standard open-world baselines (GLiNER-S/M/L, GNER-T5 and OWNER), and a 4-bit Qwen2.5-1.5B.
\item \textbf{A three-axis, seed-averaged protocol.} We report quality (AMI), latency and energy jointly rather than accuracy alone, and average every OPENER score over three embedder-retraining seeds.
\item \textbf{An explanation of why the two heads diverge, and of where quality is lost.} OPENER-ZS stays usable on English legal text (up to 49.1 gold AMI), but it does not lead the zero-shot field. It trails GLiNER-M/L and, on average, OWNER. OPENER-Sup is barely dented by the domain shift ($61.4$ mean gold AMI), and posts the best end-to-end score (35.9). Its balanced probe draws boundaries between closely related legal types that a single label-name match cannot. \rev{Splitting each end-to-end score into detection and typing shows the probe typing $78.7\%$ of detected mentions correctly, against the detector's own $64.6\%$. Even a head typing all of them correctly would stop at ${\approx}70$ AMI rather than $100$. Mentions the detector never proposes cannot be recovered downstream.} Both heads are then capped by the English detectors' recall, down to $18.7\%$ on German for GLiNER-L and $10.2\%$ for GLiNER-M. In-language \emph{detection}, not typing, is therefore the first place further adaptation should go.
\end{itemize}
